\section{Full-study protocol and exploratory diagnostics}
\label{app:transport-followup}

\subsection{Prespecified protocol, uncertainty, and provenance}

The full study fixes teacher seed 1729, development seeds 30000--30001,
tuning seeds 30100--30109, and evaluation seeds 30200--30249.
Contexts are normalized Rademacher vectors.  The feature map concatenates the
context, action indicator, and their tensor product, then normalizes it to
unit norm; its dimension is $4+5+20=29$.  The teacher has unit norm and a fixed
regular-simplex construction in the interaction block.  The learner starts at
zero, uses training ridge one, and receives only the selected action's reward.
The potential-noise table couples methods without revealing unselected rewards.
The 24 development trajectories cover three conditions and four policies at
both development seeds.  The 270 tuning trajectories cover nine optimizer
candidates, three conditions, and ten tuning seeds, using uniform-random
actions.  The nine candidates combine learning rates
$\{3\times10^{-5},10^{-4},3\times10^{-4}\}$ with $\{1,5,20\}$ full-batch steps.
All were eligible.  Candidate 008, learning rate $3\times10^{-4}$ and 20 steps,
minimized equally weighted all-action prediction MSE over the 30 tuning
seed--condition cells, with selected MSE $0.00392155$.
Evaluation consists of $3\times4\times50=600$ further trajectories.  The 294
development and tuning trajectories are not evaluation replicates.

A documented development-only numerical refinement froze Gauss--Legendre
quadrature order 2048 before tuning and evaluation.  The full development gate
compared orders 1024 and 2048 on all 24 declared trajectories; a separate
development check also compared 2048 and 4096.  No failed development check was
reclassified as a passing evaluation.  Quadrature checks, dense eigenvalue
comparisons, and confidence/optimism checks remain floating-point diagnostics.
They do not supply rigorous numerical enclosures for theorem assumptions.

The paired-bootstrap intervals resample the 50 evaluation seeds within a
condition, preserving the pairing between policies, with 10,000 resamples and
the recorded deterministic bootstrap seed.  The nine intervals in
\Cref{tab:transport-followup-outcomes} are comparisonwise rather than
multiplicity-adjusted.  Shared seeds across conditions do not make the 600
trajectories independent observations from one common population.  The teacher,
conditions, and shared optimizer are fixed in this uncertainty calculation.

The committed full-study export records the exact configuration, selected
optimizer, source inventory, development/tuning summaries, per-seed terminal
outcomes, drift summaries, and compact aggregate.  Its independent verifier
recomputes the final-horizon means and paired intervals from the per-seed
values.  The later mechanism analysis scans the preserved raw trajectories and
records their hash inventory.  The new paper tables are generated directly
from these machine artifacts by
\texttt{tools/make\_transport\_followup\_paper.py}; the accompanying
\texttt{transport\_followup\_provenance.json} binds input and output SHA-256
digests.  This provenance is separate from the original near-linear study,
whose raw trajectories were unavailable for the revision.

\subsection{Exploratory recommendations on fixed histories}

For each logged state, we evaluate all four original score rules using that
state's matrices, centers, and radii, without changing its history.
\Cref{tab:transport-followup-same-history} includes all 12 combinations of
condition and behavior policy.  Each entry averages one-step pseudo-regret
over 50,000 dependent states; the diagonal equals the actual mean cumulative
pseudo-regret divided by 1000.  Off-diagonal entries are recommendations on
states visited by the row policy.  They are not cumulative regret for an
alternative adaptive policy: a changed action would generally change future
observations, parameters, matrices, and decisions.  The analysis was specified
after the full-study outcomes were observed.

\begin{table}[t]
  \centering
  \caption{Exploratory same-history mean one-step pseudo-regret.  Rows identify
  the policy that generated the history; columns identify the score used to
  recommend an action on that state.  Bold diagonal entries are realized
  behavior-policy outcomes divided by $T$.  Off-diagonal entries do not
  reconstruct alternative adaptive runs.}
  \label{tab:transport-followup-same-history}
  % Generated by tools/make_transport_followup_paper.py; no experiments rerun.
% Input SHA-256 914013bf6436e82a50865d5ce3570b4195e7cdc2e4f25b436eb5b28a9b99ece3  review/post_full_fable/independent/analysis.json
% Input SHA-256 fd67cdc031c02876ed1e18ea0444a35348b4cb80bdd261b3fd3d75e6e920b29e  review/transport_relative_diagnostics/full/analysis.json
% Input SHA-256 9d7939364bfd7737df5daaf80d43c0673828505466ede181fe31f42fe2f98bea  review/transport_relative_diagnostics/full/per_seed.jsonl
% Input SHA-256 f28acc2d601d7524f5b30a55f396c5408a7f37a26fd36bde49a3477024bc6658  review/transport_review_followup/full-cce-full-local/aggregate_compact.json
% Input SHA-256 fb3293daacba3511ca5f3269d67e9aa0f0aca2708ad9876661ca062e197a59a4  review/transport_review_followup/full-cce-full-local/drift_and_reference_diagnostics.json
% Input SHA-256 9b6bda78c88074e46e1ff3c4253b9623bb95f467f068b6339bee9d83e4e4fce1  review/transport_review_followup/full-cce-full-local/round_drift_summaries.jsonl
% Input SHA-256 e1c4935f879df7418a2a115692113d90e50d6e50e3819cbe3ad29ce40a9afc60  review/transport_review_followup/full-cce-full-local/seed_summaries.jsonl
% Input SHA-256 b1a6e9c5d8cfe70d3c3e3bb7cbd818502c21bb3b2f276426702cb016db323ecc  review/transport_review_followup/full-cce-full-local/terminal_bound_decomposition.json
% Input SHA-256 153da33c998c3074e6fdabb9dc241157da5ff2bf0e9f401350699fa87a540496  tools/make_transport_followup_paper.py
\begin{tabular}{@{}rlrrrr@{}}
\toprule
$W$ & Logged behavior & Hessian & Endpoint & Frozen & Naive \\
\midrule
16 & Hessian & \textbf{0.185101} & 0.029032 & 0.029066 & 0.029027 \\
16 & Endpoint & 0.336227 & \textbf{0.152593} & 0.152155 & 0.151549 \\
16 & Frozen reference & 0.337046 & 0.152483 & \textbf{0.151785} & 0.151323 \\
16 & Naive current & 0.336677 & 0.153759 & 0.153626 & \textbf{0.152746} \\
\addlinespace
1 & Hessian & \textbf{0.182900} & 0.049506 & 0.059206 & 0.048660 \\
1 & Endpoint & 0.309781 & \textbf{0.173187} & 0.173884 & 0.167143 \\
1 & Frozen reference & 0.311620 & 0.174602 & \textbf{0.173195} & 0.169051 \\
1 & Naive current & 0.310986 & 0.179028 & 0.177323 & \textbf{0.173052} \\
\addlinespace
0.25 & Hessian & \textbf{0.176439} & 0.101773 & 0.144606 & 0.097656 \\
0.25 & Endpoint & 0.243915 & \textbf{0.172089} & 0.170811 & 0.164566 \\
0.25 & Frozen reference & 0.232988 & 0.175617 & \textbf{0.171876} & 0.169862 \\
0.25 & Naive current & 0.249712 & 0.180425 & 0.174521 & \textbf{0.172584} \\
\bottomrule
\end{tabular}

\end{table}

Write the primary score on a recorded state as
$m_t(a)+B_t s_t(a)+b_t$, where $b_t$ is action independent and
$B_t=\beta_t^{\rm corr}\exp(\bar D_t^Q/2)$.
Let $s_{t,(1)}\ge s_{t,(2)}$ be the largest two current widths.
The strict condition
\[
 B_t\bigl(s_{t,(1)}-s_{t,(2)}\bigr)
 >\max_a m_t(a)-\min_a m_t(a)
\]
is sufficient for the unique largest-width action to maximize the score.
Width-argmax agreement alone is weaker: prediction centers might favor the
same action.  On primary histories, agreement counts are 49,691, 49,981, and
49,991 out of 50,000 rounds, and strict-condition counts are 48,552, 49,686,
and 49,754.  The exact tie rule is the smallest action index.

The numerical scan reconstructs the recorded selected scores exactly and
finds no nonfinite coefficients or scores and no zero widths.  A separate
250-digit decimal calculation evaluates the final score arithmetic from the
exact stored binary64 centers, widths, offsets, and \emph{already rounded}
coefficient $B_t$.  It changes the primary recommendation at 16 states:
five, five, and six in the respective conditions.  Across all behavior
histories there are 65 changes among 600,000 states.  All occur at rounds
2--10; the primary differences are below 0.629 of a binary64 score ULP.
The decimal calculation does not recompute upstream matrix solves,
exponentials, coefficients, or fitted parameters.  An early changed action
could affect subsequent adaptive history; that effect has not been measured.
Likewise, equality between the full score vector and its bonus-plus-offset
part in binary64 arithmetic includes initially zero centers and is not, by
itself, evidence that the selected action was corrupted.

\Cref{tab:transport-followup-mechanism-variation} reports per-seed variation
instead of treating the pooled rounds as independent.  In every row, the
recommendation being diagnosed is the Hessian score evaluated on the row's
behavior histories.  It is the behavior policy only in Hessian rows; in other
rows its conditional recommendation can still coincide with the logged action.
Each fraction has 1000 rounds in its denominator, and the
reported median and range summarize 50 seeds.  These descriptive summaries
are not confidence intervals.

\begin{table}[t]
  \centering
  \caption{Exploratory per-seed variation in Hessian-score recommendations on
  all four behavior histories.  Entries are the median percentage of rounds
  and the minimum--maximum range across 50 seeds.  Width argmax uses the exact
  smallest-index tie rule; the strict condition is the bonus-gap inequality
  above.}
  \label{tab:transport-followup-mechanism-variation}
  % Generated by tools/make_transport_followup_paper.py; no experiments rerun.
% Input SHA-256 914013bf6436e82a50865d5ce3570b4195e7cdc2e4f25b436eb5b28a9b99ece3  review/post_full_fable/independent/analysis.json
% Input SHA-256 fd67cdc031c02876ed1e18ea0444a35348b4cb80bdd261b3fd3d75e6e920b29e  review/transport_relative_diagnostics/full/analysis.json
% Input SHA-256 9d7939364bfd7737df5daaf80d43c0673828505466ede181fe31f42fe2f98bea  review/transport_relative_diagnostics/full/per_seed.jsonl
% Input SHA-256 f28acc2d601d7524f5b30a55f396c5408a7f37a26fd36bde49a3477024bc6658  review/transport_review_followup/full-cce-full-local/aggregate_compact.json
% Input SHA-256 fb3293daacba3511ca5f3269d67e9aa0f0aca2708ad9876661ca062e197a59a4  review/transport_review_followup/full-cce-full-local/drift_and_reference_diagnostics.json
% Input SHA-256 9b6bda78c88074e46e1ff3c4253b9623bb95f467f068b6339bee9d83e4e4fce1  review/transport_review_followup/full-cce-full-local/round_drift_summaries.jsonl
% Input SHA-256 e1c4935f879df7418a2a115692113d90e50d6e50e3819cbe3ad29ce40a9afc60  review/transport_review_followup/full-cce-full-local/seed_summaries.jsonl
% Input SHA-256 b1a6e9c5d8cfe70d3c3e3bb7cbd818502c21bb3b2f276426702cb016db323ecc  review/transport_review_followup/full-cce-full-local/terminal_bound_decomposition.json
% Input SHA-256 e8e0c87f16e2bf9252b27d022a473a114c33d395cc20303d156013568e9dc836  tools/make_transport_followup_paper.py
\begin{tabular}{@{}rlrr@{}}
\toprule
$W$ & Logged behavior & \shortstack{Width argmax (\%)\\Median [min, max]} & \shortstack{Strict condition (\%)\\Median [min, max]} \\
\midrule
16 & Hessian & 99.4 [98.8, 99.9] & 97.3 [92.9, 98.6] \\
16 & Endpoint & 99.7 [99.3, 100.0] & 98.5 [97.6, 99.1] \\
16 & Frozen reference & 99.8 [99.3, 100.0] & 98.5 [97.8, 99.1] \\
16 & Naive current & 99.7 [99.3, 100.0] & 98.6 [97.4, 99.1] \\
\addlinespace
1 & Hessian & 100.0 [99.8, 100.0] & 99.4 [99.1, 99.6] \\
1 & Endpoint & 100.0 [99.7, 100.0] & 99.4 [99.1, 99.6] \\
1 & Frozen reference & 100.0 [99.7, 100.0] & 99.4 [99.2, 99.6] \\
1 & Naive current & 100.0 [99.7, 100.0] & 99.4 [99.1, 99.6] \\
\addlinespace
0.25 & Hessian & 100.0 [99.7, 100.0] & 99.5 [99.3, 99.6] \\
0.25 & Endpoint & 100.0 [99.8, 100.0] & 99.5 [99.2, 99.6] \\
0.25 & Frozen reference & 100.0 [99.7, 100.0] & 99.5 [99.3, 99.6] \\
0.25 & Naive current & 100.0 [99.7, 100.0] & 99.5 [99.3, 99.6] \\
\bottomrule
\end{tabular}

\end{table}

\subsection{Exploratory relative-certificate audit}

For unit feature and parameter bounds, \Cref{cor:tanh-relative-weights} gives
the uniform endpoint bound $D_0=4\log\cosh(1/\sqrt W)$, approximately
$0.124$, $1.735$, and $5.300$ at the three conditions.  The category-specific
sandwich $c_{-,t}\bVbar_t\preceq\bV_t\preceq c_{+,t}\bVbar_t$ gives a
symmetric coefficient $\exp(D_{{\rm cat},t}/2)$, where
$D_{{\rm cat},t}=\max\{-\log c_{-,t},\log c_{+,t}\}$, or the separate outward
coefficient $\sqrt{c_{+,t}}$ of \Cref{cor:relative-endpoint-ucb}.
The latter retains the separate return factor $c_{-,t}^{-1/2}$ in the
cumulative bound.  Neither category construction is an upper bound on the
length of the original interpolated factor path.

The post-outcome audit replays the original representation optimizer with
fixed logged actions and rewards for all 600 trajectories; it never replaces
a behavior action or generates a new observation.  It reconstructs
collection-time category weights and compares the resulting sandwich,
effective widths, and recommendations at the stored states.  Matrix, center,
and endpoint reconstruction are checked at rounds 1, 2, 10, 100, 500, and 1000;
the original float64 tolerances are retained.  These computations check an
implementation of the mathematical construction, without converting
floating-point residuals into verified certificates.  All declared checks
passed; the largest residual divided by its original tolerance was below
$0.005$.  The per-seed variation in the preceding table was reproduced by this
separate fixed-action replay.

\begin{table}[t]
  \centering
  \caption{Exploratory relative certificates on Hessian's recorded histories.
  The upper panel averages effective-width/reference-width ratios over all
  five candidate actions and 1000 rounds, then over 50 seeds per condition;
  the direct reference ratio is identically one.  The lower panel averages
  one-step pseudo-regret of recommendations on these same states.  Uniform
  uses $D_0$; the category columns use symmetric and separate outward factors.
  These are fixed-history diagnostics, not adaptive policy outcomes.}
  \label{tab:transport-followup-relative-audit}
  % Generated by tools/make_transport_followup_paper.py; no experiments rerun.
% Input SHA-256 914013bf6436e82a50865d5ce3570b4195e7cdc2e4f25b436eb5b28a9b99ece3  review/post_full_fable/independent/analysis.json
% Input SHA-256 fd67cdc031c02876ed1e18ea0444a35348b4cb80bdd261b3fd3d75e6e920b29e  review/transport_relative_diagnostics/full/analysis.json
% Input SHA-256 9d7939364bfd7737df5daaf80d43c0673828505466ede181fe31f42fe2f98bea  review/transport_relative_diagnostics/full/per_seed.jsonl
% Input SHA-256 f28acc2d601d7524f5b30a55f396c5408a7f37a26fd36bde49a3477024bc6658  review/transport_review_followup/full-cce-full-local/aggregate_compact.json
% Input SHA-256 fb3293daacba3511ca5f3269d67e9aa0f0aca2708ad9876661ca062e197a59a4  review/transport_review_followup/full-cce-full-local/drift_and_reference_diagnostics.json
% Input SHA-256 9b6bda78c88074e46e1ff3c4253b9623bb95f467f068b6339bee9d83e4e4fce1  review/transport_review_followup/full-cce-full-local/round_drift_summaries.jsonl
% Input SHA-256 e1c4935f879df7418a2a115692113d90e50d6e50e3819cbe3ad29ce40a9afc60  review/transport_review_followup/full-cce-full-local/seed_summaries.jsonl
% Input SHA-256 b1a6e9c5d8cfe70d3c3e3bb7cbd818502c21bb3b2f276426702cb016db323ecc  review/transport_review_followup/full-cce-full-local/terminal_bound_decomposition.json
% Input SHA-256 e8e0c87f16e2bf9252b27d022a473a114c33d395cc20303d156013568e9dc836  tools/make_transport_followup_paper.py
\begin{tabular}{@{}rrrrr@{}}
\toprule
$W$ & Dense endpoint & Uniform & Category symmetric & Category asymmetric \\
\midrule
\multicolumn{5}{c}{Mean effective width / reference width} \\
\addlinespace
16 & 1.0016 & 1.0638 & 1.0033 & 1.0028 \\
1 & 1.0264 & 2.3819 & 1.0555 & 1.0427 \\
0.25 & 1.1213 & 14.2984 & 1.3313 & 1.1268 \\
\addlinespace
\multicolumn{5}{c}{Mean one-step pseudo-regret on the same fixed states} \\
\addlinespace
16 & 0.029032 & 0.029524 & 0.029052 & 0.029047 \\
1 & 0.049506 & 0.086686 & 0.050706 & 0.050322 \\
0.25 & 0.101773 & 0.168958 & 0.109173 & 0.102158 \\
\bottomrule
\end{tabular}

\end{table}

\Cref{tab:transport-followup-relative-audit} uses the primary policy's
histories; the machine artifact also reports all four behavior-history groups
and per-seed summaries.  Direct frozen-reference recommendations on these
states appear in \Cref{tab:transport-followup-same-history}.  The audit also
records timings for category-ratio kernels with and without the 80 feature
projections, and for dense endpoint evaluation with preassembled matrices.
These are fixed-state microbenchmarks under parallel analysis load; they
exclude matrix construction, confidence estimation, and the policy loop.
Over 36 states (one seed, three checkpoints, four behaviors, and three
conditions), the median of the per-state median kernel times was
$27.7\,\mu\mathrm{s}$ for category ratios with precomputed scales,
$34.8\,\mu\mathrm{s}$ including 80 projections, and
$213.3\,\mu\mathrm{s}$ for dense endpoint evaluation.  Each state's timing
uses five repeated batches; these timings are descriptive observations on
one host, not a statistical comparison of end-to-end algorithms.
The category construction requires additional accumulated collection weights
and still uses matrix solves and the frozen confidence construction.  No
end-to-end speedup or adaptive regret improvement has been established.
